\section{Worked finite compilations}

The examples in this section are small enough to inspect manually, but each exercises a distinct generic layer.  Their exact host-level computations are recorded in \path{examples/demo_certificate.json}.

\subsection[A bounded forall-exists strategy]{A bounded \(\forall\exists\) strategy}

Consider
\[
  \forall x<6\;\exists y<4\quad (x+y)\bmod 3=0.
\]
The matrix is decidable by the checked division/remainder arithmetic.  For each \(x\), bounded least search returns the least \(y<4\) satisfying the congruence.  The resulting table is
\[
  \sigma=[0,2,1,0,2,1].
\]
Indeed,
\[
0+0,\ 1+2,\ 2+1,\ 3+0,\ 4+2,\ 5+1
\]
are all divisible by three.  The table verifier checks six bounds and six remainder equations.  By finite strategy soundness, the verifier implies the quantified formula.

There are several possible public endpoints:

\begin{itemize}
\item prove the bounded formula and keep \(\sigma\) private;
\item return \(\exists\sigma\,\mathsf{Verifies}(\sigma)\);
\item return the least strategy table by canonical search;
\item export only a stable statistic of the table.
\end{itemize}

The choice among them is a theorem-design decision, not an implementation accident.

\subsection{A cyclic precision network}

Take caps \((5,6,5)\) and demand transformer
\[
F(d_0,d_1,d_2)=
\bigl(\min(5,d_1+1),\min(6,d_2+2),\min(5,d_0)\bigr).
\]
It is monotone.  Starting at zero and applying \(H(d)=d\vee F(d)\) gives
\[
\begin{array}{c|c|c}
 t & d_t & \rank(d_t)\\
\hline
0&(0,0,0)&0\\
1&(1,2,0)&3\\
2&(3,2,1)&6\\
3&(3,3,3)&9\\
4&(4,5,3)&12\\
5&(5,5,4)&14\\
6&(5,6,5)&16\\
7&(5,6,5)&16.
\end{array}
\]
The total cap rank is \(5+6+5=16\), and the fixed point is reached after six strict iterations.  The example demonstrates why one reverse pass is insufficient and why the rank budget gives a uniform stopping bound.

In an analytic graph, the coordinates could represent input precision, range-bound precision, and a denominator-separation precision.  The fixed point is computed before the forward numerical trace, so the trace invariant mentions only concrete resource terms.

\subsection{Viability versus greedy choice}

Let
\[
A_0=\{0,1\},\qquad A_1=\{0,1\},\qquad A_2=\{0\}.
\]
All nodes are locally valid.  The compatibility edges are
\[
  0\to0,\quad 1\to1
  \qquad\text{between }A_0,A_1,
\]
and only
\[
  1\to0
  \qquad\text{between }A_1,A_2.
\]
Greedy least local choice starts with \(0\), moves to \(0\), and fails.  Backward viability yields
\[
  V_2=\{0\},\qquad V_1=\{1\},\qquad V_0=\{1\}.
\]
The least viable path is \((1,1,0)\).  The example is deliberately included in the formalization acceptance gates because it catches a subtle but common incorrect selector.

\subsection{Metric observation on a finite tree}

Take a rooted tree with edges \(r\to a\), \(r\to b\), and \(a\to c\).  Suppose certified edge errors are
\[
 e_a=2,\qquad e_b=1,\qquad e_c=3.
\]
Metric tree descent gives root-distance bounds
\[
  E_r=0,
  \quad E_a=2,
  \quad E_b=1,
  \quad E_c=e_a+e_c=5.
\]
The public result can expose the root observation and these explicit error radii.  It need not choose globally compatible representatives on every cycle of a larger cover unless the downstream theorem actually requires them.

\subsection{Height-bounded Diophantine enumeration}

Let \(P(x)\) be the reflected equation predicate for a fixed Diophantine problem under a natural code.  Assume a residual formula
\[
  \mathsf{HB}(H):=\forall x\,[P(x)\to x<H].
\]
The compiler filters \(0,\ldots,H-1\), producing a list \(L_H\).  Bounded reflection proves:
\[
  x\in L_H\leftrightarrow x<H\wedge P(x).
\]
Under \(\mathsf{HB}(H)\), this becomes \(x\in L_H\leftrightarrow P(x)\).  The result is the conditional object theorem
\[
  \mathsf{HB}(H)\vdash_{\PRAAx}\mathsf{Complete}(L_H,P).
\]

Nothing in the finite compiler claims to produce \(H\).  A theorem-specific descent, Runge estimate, Baker bound, or other effective argument can later discharge the residual.

\subsection{A quantitative bisection contract}

Let \(f\) be a fixed primitive-recursive rational function on a rational interval \([a,b]\).  Suppose the residual context supplies:

\begin{enumerate}
\item opposite endpoint signs;
\item a monotonicity certificate on the interval;
\item a denominator/apartness guard for every rational operation;
\item a primitive-recursive bound \(N(p)\) such that \((b-a)/2^{N(p)}\le2^{-p}\).
\end{enumerate}

A trace stores endpoints \((\ell_t,r_t)\).  At each step it computes the midpoint and uses the reflected sign bit to retain a bracketing half.  The quantifier-free invariant states:
\[
 a\le\ell_t\le r_t\le b,
 \quad f(\ell_t)\le0\le f(r_t),
 \quad 2^t(r_t-\ell_t)=b-a.
\]
After \(N(p)\) steps, return the midpoint and radius \(2^{-p}\).  Monotonicity proves every exact root in the bracket lies in the output ball; a separation/strict-monotonicity modulus can strengthen this to uniqueness.  The theorem constructs a certified finite approximation without postulating a completed real root.

\subsection{End-to-end composition}

A representative SCOPE compilation has the following order:

\begin{enumerate}
\item specialize the theorem graph to query \(q\);
\item compute capped cyclic demand closure;
\item decode or produce a bounded master certificate;
\item reflect all bounded validity conditions;
\item compute viability for nonunique paths;
\item run fixed forward traces at the closed resources;
\item project through an exact or metric observer;
\item eliminate private certificates;
\item retain only reachable producer residuals;
\item emit the final \(\Prf\,\PRAAx\) derivation.
\end{enumerate}

The demo JSON records finite evidence for steps 2, 5, 6 in strategy form, and 7 in metric-tree form.  It is executable evidence for the algorithms, not a substitute for the object derivations described above.
